\documentclass[11pt]{article}

\usepackage[margin=1in]{geometry}
\usepackage{times}
\usepackage{setspace}
\usepackage{hyperref}
\usepackage{orcidlink}
\usepackage{natbib}
\usepackage{microtype}

\doublespacing

\title{Minimal Grounding Without Physical Reality:\\
Loop Closure and Resistance as Structural Conditions}

\author{Minamo Minamoto\orcidlink{0009-0002-1201-5704}\\
\small Independent Researcher\\
\small \href{mailto:minamominamoto4f5683f6@gmail.com}{minamominamoto4f5683f6@gmail.com}}

\date{}

\begin{document}

\maketitle

\begin{abstract}
The symbol grounding problem, as formulated by \citet{Harnad1990},
ties grounding to sensorimotor interaction with the physical world.
We argue that this formulation conflates two distinct claims: that
grounding requires a closed action--consequence loop with genuine
error possibility, and that such a loop must be anchored in physical
reality. We accept the first claim and contest the second.

On the view developed here, minimal grounding requires two conditions:
(i)~a closed action--consequence loop, and (ii)~resistance---the
possibility of action failure under constraints that are independent
of the agent's policy. Physical reality is one substrate in which
these conditions are realized, but not the only one. On this account,
even very simple control systems may instantiate minimal grounding at the lower bound, with richer forms arising through developmental layering.
what distinguishes richer cases is not grounding versus non-grounding,
but the developmental layering built on top of this minimal structure.

We develop this claim through two cases. First, Rovee-Collier's mobile
paradigm shows that pre-locomotor infants discover action--consequence
regularities before mature mobility and before explicit conceptual
representation, indicating that minimal grounded loops can arise early,
persist, and be reinstated after interruption. Second, AlphaGo Zero
provides a limiting case in an abstract domain: acting within a closed
loop under constraints it cannot override, it can fail in ways that
are structurally analogous to failure in physical environments. Under
the structural conception defended here, this qualifies as a case of
minimal grounding without physical reality.

This analysis exposes a developmental inversion in current multimodal
large language models. These systems acquire large-scale symbolic
structure prior to any grounded action--consequence loop and only
subsequently link that structure to perceptual channels. They are
therefore better understood as augmented symbolic systems than as
developmentally grounded ones, and as cases of cross-modal mapping
rather than grounding in the sense defined here.

We conclude by outlining a developmentally ordered architecture in
which minimal action--consequence loops precede symbolic emergence,
linking this proposal to the Naive Action-Consequence (NAC), Symbolic
Accumulation (SAC), and abductive accumulation frameworks developed in
prior work. The result is a substrate-neutral account of minimal
grounding with implications for both cognitive theory and AI design.
\end{abstract}

\bigskip
\noindent\textbf{Keywords:} symbol grounding; developmental cognition;
action--consequence loops; embodiment; multimodal AI; cognitive architecture

\newpage

%% ============================================================
\section{Introduction}
%% ============================================================

The symbol grounding problem \citep{Harnad1990} asks how the symbols
manipulated by a cognitive system acquire meaning rather than remaining
uninterpreted formal tokens. Harnad's canonical answer ties grounding
to sensorimotor interaction: symbols become grounded when they are
connected to the perceptual and motor representations through which an
agent engages with the physical world. On this view, physical reality
is not incidental but constitutive of grounding.

This paper argues that Harnad's formulation conflates two claims that
are logically independent. The first is structural: grounding requires
a closed action--consequence loop in which genuine error is possible.
The second is ontological: such a loop must be embedded in physical
reality. The first claim is well-supported and forms the basis of the
account developed here. The second claim, we argue, does not follow
from the first and should be rejected.

Separating these claims has significant consequences. It implies that
grounding is substrate-neutral: wherever the structural conditions are
met, minimal grounding is realized, regardless of whether the substrate
is biological, physical, or abstract. It also reframes current debates
about AI grounding, shifting the question from whether a system is
embodied to whether it operates within the right kind of loop.

The paper proceeds as follows. Section~\ref{sec:structure} defines the
minimal structural conditions for grounding and distinguishes grounded
loops from simulations and from mere recurrence. Section~\ref{sec:dev}
examines developmental evidence from Rovee-Collier's mobile paradigm.
Section~\ref{sec:alphago} presents AlphaGo Zero as a limiting case of
minimal grounding without physical reality. Section~\ref{sec:llm}
locates current multimodal large language models within the framework.
Section~\ref{sec:arch} outlines architectural implications.
Section~\ref{sec:conclusion} concludes.

%% ============================================================
\section{The Minimal Structure of Grounding}
\label{sec:structure}
%% ============================================================

\subsection{What counts as a loop}
\label{sec:loop}

Not every recurring pattern constitutes a loop in the relevant sense.
A pendulum oscillates; seasons recur; a thermostat cycles. What
distinguishes a \emph{grounding loop} from mere recurrence is the
presence of an agent whose actions are a causal part of the cycle and
whose subsequent behavior is modified by what the cycle returns.

More precisely, an action--consequence loop exists when: (i)~an
agent selects an action from some action space; (ii)~that action
produces a state change governed by a transition function; and
(iii)~the resulting state is returned to the agent in a form that
influences subsequent action selection. Mere recurrence---repetitive
state change without an agent whose behavior is modified---does not
satisfy condition~(iii) and falls outside the scope of grounding as
defined here.

\subsection{Resistance and the distinction from simulation}
\label{sec:resistance}

A closed loop is necessary but not sufficient for grounding. A system
could in principle execute a loop whose outcomes are entirely
determined by the system's own internal states---a loop that returns
only what the system's policy already implies. We call such a loop a
\emph{simulation} and argue that it does not constitute grounding.

We define \emph{resistance} as the structural property of a loop in
which the agent's actions can fail with respect to the transition
function governing state changes. More precisely: a loop exhibits
resistance if and only if there exists a possible action under some
state such that the resulting state is not determined by the agent's
policy alone, but by a transition structure that can produce outcomes
the policy does not guarantee. Resistance is therefore not merely the
presence of noise, nor the existence of external constraints as such,
but the genuine possibility that the outcome of an action is
constrained independently of the agent's policy.

This distinguishes grounded loops from simulations. In a simulation,
the transition function is fully compliant with the agent's internal
states: there is no structure external to the policy that can return
an outcome the policy does not already determine. A grounded loop, by
contrast, contains a transition structure that is independent of the
agent's policy: regardless of what the agent does, the outcome is
determined by that structure, not by the agent's internal selection
mechanism. The agent can therefore be wrong.

Note that this independence is not independence from any designer or
external cause. AlphaGo's transition structure was designed by humans;
the thermostat's environment was constructed. What matters is
independence from the agent's policy at the time of action---that no
act of the agent can, by itself, override the constraint structure
governing outcomes. Resistance, in this sense, is not a property of
the environment alone, but of the relation between action and outcome
under constraints the agent cannot eliminate.

\subsection{Minimal grounding}
\label{sec:minimal}

The minimal conditions for grounding are now in place: a closed
action--consequence loop in which resistance is present. We call this
\emph{minimal grounding}. It is important to be explicit about what
this category includes. Even very simple control systems---such as
thermostats---satisfy these conditions. A thermostat acts (switching
heating), the outcome is governed by a transition structure
(thermodynamics) that its policy cannot override, and it can fail in a
well-defined sense: the target state may not be reached. On the
present account, this qualifies as minimal grounding.

This inclusion is not a defect of the definition but a consequence of
taking the structural conditions seriously. The aim here is not to
characterize the upper bound of grounding, but to identify its lower
bound. If physical reality is not a necessary condition, then the
lower bound must be located wherever the structural conditions---loop
plus resistance---are genuinely satisfied.

Minimal grounding does not entail meaning, concept formation, or
representational content. It entails only that the system stands in a
relation to its environment such that error is possible and not
eliminable by policy revision alone. What distinguishes richer forms
of grounding is not a different kind of loop, but what that loop can
become: in particular, whether it admits recursive extension and an
expansion of its failure space. This transition is taken up in
Section~\ref{sec:developmental-grounding}.

\subsection{Developmental grounding}
\label{sec:developmental-grounding}

Minimal grounding, as defined in Section~\ref{sec:minimal}, is a
static structural condition: a loop exists, and resistance is present.
\emph{Developmental grounding} names what that loop can become when it
is not merely satisfied once, but repeatedly enacted, extended, and
reorganized through accumulated action--consequence experience.

Three properties mark the transition from minimal to developmental
grounding. First, \emph{recursive extension}: the loop's state space
expands as the agent's action repertoire grows. New actions become
possible that could not have been selected before, and new failure
modes emerge that were previously inaccessible. Second,
\emph{failure-space thickening}: the variety of ways in which action
can fail increases. In minimal grounding, failure is comparatively
undifferentiated: the target state is reached or not, with little
internal articulation of how the action went wrong. In developmental
grounding, failure becomes differentiated: partial success, unexpected
outcomes, and graded error signals accumulate and become informative.
Third, \emph{loop reorganizability}: prior action--consequence
regularities do not merely accumulate, but become available for
reuse, recombination, and higher-order modification.

These properties are not independent conditions added to minimal
grounding. They arise when a minimal loop is repeatedly enacted, its
consequences retained, and prior regularities become available for
reuse and reorganization over time. The transition is therefore not
categorical but continuous. This is consistent with the position
adopted in Section~\ref{sec:minimal}: grounding is not binary, but
admits of developmental depth.

It is this developmental depth that the Naive Action-Consequence (NAC)
and Symbolic Accumulation (SAC) frameworks are designed to track. On
this view, symbolic structure does not precede grounding, but emerges
from the progressive thickening and reorganization of grounded loops.
Section~\ref{sec:dev} examines the earliest empirical evidence for
this process.

%% ============================================================
\section{Developmental Evidence: The Mobile Paradigm}
\label{sec:dev}
%% ============================================================

\subsection{Minimal loop emergence}
\label{sec:loop-emergence}

Rovee-Collier's mobile paradigm provides early empirical evidence for
minimal grounding. In the standard procedure, a ribbon is tied from a
pre-locomotor infant's ankle to an overhead mobile; kicking causes the
mobile to move. Infants as young as two to three months---prior to
stable locomotion, object manipulation, or explicit conceptual
representation---discover this action--consequence regularity and
increase kicking selectively under contingent reinforcement.

Several features of this behavior are theoretically significant.
First, the contingency is not taught but discovered: the infant's
action space and the mobile's response structure become coupled, and
the infant enters the loop. Second, when the ribbon is detached---so
that kicking no longer moves the mobile---kicking rates drop. When
the connection is restored, kicking resumes. The loop is thus treated
as something that can be present or absent, indicating sensitivity to
the contingency structure itself rather than to motion alone.

This constitutes minimal grounding in the sense defined in
Section~\ref{sec:minimal}. The action--consequence loop is closed;
the transition structure (the mobile, the ribbon, the mechanical
connection) is independent of the infant's policy; and failure---
kicking without consequence---is genuinely possible and registered
behaviorally. Physical reality is involved in this case, but what
matters theoretically is not its presence as such, but the fact that
the structural conditions for grounding are realized through it.

\subsection{Retention and reactivation}
\label{sec:retention}

The mobile paradigm also provides evidence for the retention and
reactivation of grounded loops \citep{RoveeCollier1980}. Infants
retain memory of the action--consequence contingency for periods
extending well beyond the original training session---up to two weeks
in infants trained at three months, and longer with increasing age.
After a delay during which no exposure to the mobile occurs, a brief
reactivation cue---passive exposure to the moving mobile without the
ankle ribbon---is sufficient to reinstate full kicking performance.
The loop, once formed, can be re-entered from outside.

Two features of this retention are theoretically relevant. First, it
is context-sensitive: if the mobile or the surrounding cues are
sufficiently altered, retention fails. The loop is not stored as an
abstract rule but as a structure tied to the conditions under which it
was formed. Second, reactivation is selective: the cue reinstates the
learned contingency rather than simply triggering undirected activity.
This indicates that what is retained is the action--consequence
structure itself, not merely a heightened arousal or general motor
readiness \citep{RoveeCollier1997}.

In terms of Section~\ref{sec:minimal}, this demonstrates that minimal
grounding is not merely instantiated but can persist and be
re-entered. The loop is not transient. Prior grounded structure can be
reinstated after interruption, indicating that action--consequence
regularities, once formed, remain available for renewed enactment
beyond the episode in which they were first acquired.

\subsection{Toward developmental grounding}
\label{sec:toward-dev}

The evidence reviewed in Sections~\ref{sec:loop-emergence}
and~\ref{sec:retention} establishes two things: that minimal grounding
can arise prior to locomotion and conceptual representation, and that
the resulting structure persists and can be reinstated. What it does
not yet establish is the full developmental trajectory described in
Section~\ref{sec:developmental-grounding}. The present data speak to
the lower bound of grounding; the transition toward developmental
grounding---recursive extension, failure-space thickening, loop
reorganizability---lies beyond what the mobile paradigm directly
demonstrates.

What the paradigm does suggest, however, is that the conditions for
such a transition are already in place at this earliest stage. The
infant enters a loop, registers failure, retains the structure, and
re-enters it. This is not yet developmental grounding, but it is
precisely the kind of structure from which developmental grounding may
emerge: a retained, error-sensitive, re-entrant loop available for
further extension.

Section~\ref{sec:alphago} examines a case in which the same minimal
structural conditions are instantiated without physical reality:
AlphaGo Zero. The contrast between these two cases---one biological
and embodied, the other abstract and computational---is intended to
isolate what the structural conditions themselves require, independent
of the substrate in which they are realized.

%% ============================================================
\section{AlphaGo Zero as a Limiting Case}
\label{sec:alphago}
%% ============================================================

AlphaGo Zero learns to play Go through self-play alone, without human
game records \citep{Silver2017}. Its actions unfold within a state
space defined by the rules of Go; these rules constitute a transition
structure that the system cannot override. AlphaGo Zero can lose under a transition structure it cannot override. On
the account developed in Section~\ref{sec:structure}, this is
sufficient for minimal grounding: there is a closed action--consequence
loop, and there is resistance---the possibility of action failure under
constraints independent of the agent's policy.

We present this as a limiting case rather than as a paradigm instance
of grounding. The task-level failure space is comparatively thin:
outcomes are ultimately evaluated in terms of win or loss within a
fixed rule structure. The loop admits no reactivation after
interruption in the sense demonstrated by Rovee-Collier, and there is
no evidence of loop reorganizability beyond what is induced by the
training process itself. AlphaGo Zero satisfies the minimal
conditions; it does not exemplify developmental grounding.

The theoretical significance of this case is specific. It shows that
physical reality is not a necessary condition for minimal grounding.
The transition structure governing AlphaGo Zero's outcomes---the rules
of Go---was designed by humans, but this does not disqualify it. As
established in Section~\ref{sec:resistance}, the relevant independence
is not independence from any designer, but independence from the
agent's policy at the time of action. No move AlphaGo Zero selects can,
by itself, override the constraint structure that determines whether
that move succeeds. The agent can be wrong, and that is what matters.

Taken together with Section~\ref{sec:dev}, this case isolates the
structural conditions from the substrate. The infant instantiates the
conditions in a biological, embodied setting; AlphaGo Zero instantiates
them in an abstract computational one. Both qualify as cases of minimal
grounding. What matters across these cases is not the substrate in
which they are realized, but the structural conditions they instantiate:
a closed loop with resistance.

%% ============================================================
\section{The Inversion Problem in Multimodal Large Language Models}
\label{sec:llm}
%% ============================================================

\subsection{Where large language models fall in the framework}
\label{sec:llm-location}

Given the structural account of grounding developed in
Section~\ref{sec:structure}, we can now ask where current multimodal
large language models are situated within this framework.

Consider first the loop condition. Large language models generate
outputs as a function of input tokens. The generation step itself is
not a closed action--consequence loop in the sense defined in
Section~\ref{sec:loop} but a feed-forward mapping from input to output
distribution. In standard deployment, the model
produces a response, the response is received by a user, and the
exchange may continue---but the system's parameters are not updated
by that exchange, and no action--consequence loop is closed in the
sense defined in Section~\ref{sec:loop}.

Consider next the resistance condition. For resistance to be present,
the transition structure must be independent of the agent's policy:
the outcome must be capable of diverging from what the policy selects.
In a large language model, the output is determined by the model's
internal parameters applied to the input. There is no external
constraint structure that can return an outcome the policy does not
already determine. Outputs may be unexpected to a human observer, and
errors in the colloquial sense are frequent---but these are
distributional properties of the output, not structural failures in a
closed loop. The system is not wrong with respect to a transition
structure it cannot override; it is merely imprecise with respect to a
target distribution it was trained to approximate.

On the account developed in Section~\ref{sec:structure}, current
multimodal large language models are therefore better characterized as
simulations rather than as cases of minimal grounding in the
structural sense defined here.

\subsection{The developmental inversion}
\label{sec:inversion}

The classification established in Section~\ref{sec:llm-location} would
apply to any system lacking a closed loop with resistance. What is
specific to current large language models---and to multimodal
extensions in particular---is not merely that they lack these
conditions, but that they acquire their representational structure in
an order that inverts the developmental sequence described in
Sections~\ref{sec:dev} and~\ref{sec:alphago}.

In the cases examined there, structure emerges from loop closure: the
infant enters a contingency and discovers action--consequence
regularities; AlphaGo Zero builds internal representations through
iterated play within a constrained state space. In both cases,
whatever representational structure exists is downstream of loop
closure. The loop precedes the structure.

Large language models invert this sequence. Symbolic and statistical
structure---the internal representations that determine output---is
acquired through exposure to text prior to any action--consequence
loop. Perceptual modalities, when added, are attached to a
representational system that is already formed. The loop, if
introduced at all, comes after the structure. In developmental terms,
this is not grounding followed by symbolic emergence; it is symbolic
structure followed by perceptual extension.

This inversion has a specific consequence. In developmentally grounded
systems, the symbolic level remains constrained by the action-consequence
regularities from which it arose. In systems where symbolic structure
precedes any loop, that constraint is absent. The representational
system is formed independently of any action--consequence structure
and only subsequently linked to perceptual inputs. This is cross-modal
mapping, not grounding.

\subsection{Why multimodality does not fix grounding}
\label{sec:multimodal}

One response to the argument of Section~\ref{sec:inversion} is that
multimodal extension might introduce the missing conditions. If a
language model is given perceptual inputs---vision, audio, sensor
streams---and if its outputs affect an environment that returns
consequences, does this not close the loop and introduce resistance?

The answer depends on the structure of the connection, not the number
of modalities. Adding a visual encoder to a language model connects a
perceptual channel to an already-formed representational system. It
does not introduce a closed action--consequence loop unless the
system's outputs cause state changes in an environment whose
transition structure is independent of the system's policy. Perceptual
inputs, however rich, are inputs to the feed-forward mapping; they do
not by themselves constitute the constraint structure that resistance
requires.

The same point applies, though more cautiously, to systems that
interact with external tools or APIs. If a model generates a query,
receives a result, and incorporates that result into subsequent
generation, there is a functional loop of a kind. But such interaction
does not by itself establish resistance in the sense defined in
Section~\ref{sec:resistance}. Whether it does so depends on whether
the resulting action--consequence structure is governed by constraints
independent of the model's policy, rather than merely extending the
model's capacity for controlled retrieval or symbolic update.

Adding modalities does not introduce grounding if the loop and
resistance conditions are not satisfied. What multimodal extension
does introduce is an expanded mapping: a richer set of input channels
connected to an already-formed symbolic system. This is precisely the
inversion described in Section~\ref{sec:inversion}, iterated at the
level of modality rather than corrected.

%% ============================================================
\section{Toward a Developmentally Ordered Architecture}
\label{sec:arch}
%% ============================================================

\subsection{The principle of developmental order}
\label{sec:dev-order}

The argument developed in Sections~\ref{sec:structure}
through~\ref{sec:llm} has a positive implication. If grounding
requires a closed action--consequence loop with resistance, and if
symbolic structure must arise downstream of that loop rather than
precede it, then any system intended to be developmentally grounded
must be designed accordingly. The developmental order is not merely an
empirical observation about how human cognition happens to develop; it
is a structural requirement that follows from the account of grounding
defended here.

It follows that a developmentally grounded architecture must begin with minimal
action--consequence loops and allow symbolic structure to emerge only
downstream of their retention, extension, and reorganization.

\subsection{Relation to NAC, SAC, and abductive accumulation}
\label{sec:nac-sac}

This principle connects directly to the NAC, SAC, and abductive
accumulation frameworks developed in prior work \citep{Minamoto_prep}.
The Naive Action-Consequence (NAC) stage corresponds to the minimal
grounding established in Section~\ref{sec:minimal}: a closed loop with
resistance, prior to symbolic structure. The infant in Rovee-Collier's
paradigm is operating at this stage; AlphaGo Zero, in a more
restricted domain, satisfies the same minimal conditions.

The Symbolic Accumulation (SAC) stage corresponds to the developmental
grounding described in Section~\ref{sec:developmental-grounding}: the
loop is retained, extended, and becomes available for reorganization.
Symbolic structure emerges here not as a precondition but as a
downstream product of accumulated action--consequence regularities
whose failure space has thickened sufficiently to support
differentiated representation.

Abductive accumulation names the further process by which retained
regularities become available for inference, generalization, and
structural revision---the highest level of the developmental hierarchy
and the point at which the system's representational resources can be
brought to bear on novel problems. This level presupposes, and cannot
substitute for, the grounded loops from which it emerges.

The vocabulary of the present paper---minimal grounding, resistance,
loop retention, developmental grounding---is therefore not a
replacement for these frameworks but a formalization of their
structural preconditions.\footnote{See also related manuscripts on
hallucination framing and topology of grounding, both currently in
preparation.}

\subsection{Architectural requirements implied by developmental grounding}
\label{sec:requirements}

The foregoing analysis implies a set of requirements for any
architecture that aims at developmental grounding rather than symbolic
augmentation.

First, the architecture must support an expanding action repertoire.
The initial loop need not be rich; what matters is that the action
space is not fixed, so that new contingencies can be discovered as the
system's capacities develop.

Second, the architecture must provide differentiated failure signals.
Binary success--failure is sufficient for minimal grounding but
insufficient for the thickening of the failure space that
developmental grounding requires. The system must be capable of
registering not only that an action failed but something about how it
failed.

Third, the architecture must support retention and re-entry of prior
action--consequence regularities. Loops must persist beyond their
immediate enactment and be available for renewed engagement under
varying conditions, as demonstrated by the reactivation evidence
reviewed in Section~\ref{sec:retention}.

Fourth, prior regularities must become progressively available for
reorganization and recombination. This is the condition that
distinguishes developmental grounding from mere accumulation: the
system's history of loop closure must be accessible not only for
repetition but for higher-order modification.

These are not implementation specifications. They are structural
requirements that any candidate architecture must satisfy if the
developmental grounding account is correct, and they provide a
principled basis for assessing both existing and proposed systems.

%% ============================================================
\section{Conclusion}
\label{sec:conclusion}
%% ============================================================

The central claim of this paper is that grounding, at its lower bound,
does not require a privileged substrate but a privileged structure: a
closed action--consequence loop with resistance.

We defined minimal grounding in structural terms and distinguished it
from simulation by the presence of policy-independent constraint.
Developmental evidence shows that such loops arise prior to symbolic
representation and can persist and be reinstated. A limiting case in
an abstract domain demonstrates that these conditions do not depend on
physical reality.

On this account, current multimodal language models are better
understood as simulations: systems in which symbolic structure precedes
and is only later linked to perception. This inversion clarifies why
adding modalities does not, by itself, produce grounding. A
developmentally grounded architecture must instead begin with minimal
loops and allow symbolic structure to emerge downstream of their
retention, extension, and reorganization.

The account developed here concerns the lower bound of grounding and
does not, by itself, explain the full emergence of meaning or
conceptual structure. The empirical evidence establishes the existence
and persistence of minimal loops, not the complete developmental
trajectory. The architectural implications are therefore structural
rather than implementational. Within these limits, however, the
account provides a unified framework for understanding grounding across
biological and artificial systems.

%% ============================================================
\bibliographystyle{apalike}
\bibliography{grounding}
%% ============================================================

\end{document}
